\section{级数}
\label{级数}
\subsection{复数项级数}
\label{复数项级数}
\subsubsection{复数列的极限}
\label{复数列的极限}
设 \(\{\alpha_n\} (n=1,2,\cdots)\) 为一复数列，其中
\(\alpha_n = a_n + ib_n\)，又设 \(\alpha = a + ib\)
为一确定的复数。如果任意给定 \(\varepsilon > 0\)，相应地都能找到一个正数
\(N(\varepsilon)\)，使 \(|\alpha_n - \alpha| < \varepsilon\) 在
\(n > N\) 时成立，那么 \(\alpha\) 称为复数列 \(\{\alpha_n\}\) 当
\(n \to \infty\) 时的极限，记作

\[
\lim_{n \to \infty} \alpha_n = \alpha
\]

此时也称复数列 \(\{\alpha_n\}\) 收敛于 \(\alpha\)。

\begin{tcolorbox}[title={Note},colback=red!5!white
,colframe=red!75!black
,colbacktitle=yellow!50!red
,coltitle=red!25!black, fonttitle=\bfseries,subtitle style={boxrule=0.4pt, colback=yellow!50!red!25!white}]
复数列收敛的充要条件是
\[
\boxed{\lim_{n\to \infty}a_n=a,\quad \lim_{n\to \infty}b_n=b.}
\]
也就是说，复数列的敛散性可以转化为两个实数列的敛散性。

\end{tcolorbox}
\subsubsection{级数的概念}
\label{级数的概念}
\paragraph{级数定义}
\label{级数定义}
设 \(\{\alpha_n\} = \{a_n + ib_n\} (n=1,2,\cdots)\) 为一复数列，表达式

\[
\sum_{n=1}^{\infty} \alpha_n = \alpha_1 + \alpha_2 + \cdots + \alpha_n + \cdots
\]

称为\textbf{复数项无穷级数}。其最前面 \(n\) 项的和

\[
s_n = \alpha_1 + \alpha_2 + \cdots + \alpha_n
\]

称为级数的\textbf{部分和}。
\paragraph{收敛与发散}
\label{收敛与发散}
如果部分和数列 \(\{s_{n}\}\) 收敛，那么级数
\(\sum\limits_{n=1}^{\infty}\alpha_{n}\) 收敛，并且极限
\(\lim\limits_{n\to\infty}s_{n}=s\) 称为级数的和。

如果部分和数列 \(\{s_{n}\}\) 不收敛，那么级数
\(\sum\limits_{n=1}^{\infty}\alpha_{n}\) 发散。
\paragraph{收敛的条件}
\label{收敛的条件}
级数 \(\sum\limits_{n=1}^{\infty} \alpha_n = \sum\limits_{n=1}^{\infty} (a_n + i b_n)\)
收敛的充要条件是：级数 \(\sum\limits_{n=1}^{\infty} a_n\) 和 \(\sum\limits_{n=1}^{\infty} b_n\) 都收敛。

\begin{tcolorbox}[title={必要条件},colback=red!5!white
,colframe=red!75!black
,colbacktitle=yellow!50!red
,coltitle=red!25!black, fonttitle=\bfseries,subtitle style={boxrule=0.4pt, colback=yellow!50!red!25!white}]
因为实数项级数\(\sum\limits_{n=1}^{\infty} a_{n}\)和\(\sum\limits_{n=1}^{\infty} b_{n}\)收敛的必要条件是\(\lim\limits_{n\to\infty} a_{n}=0\)和\(\lim\limits_{n\to\infty}b_{n}=0\)。

所以复数项级数\(\sum\limits_{n=1}^{\infty} \alpha_{n}\)收敛的必要条件是

\[
\boxed{\lim\limits_{n\to\infty} \alpha_{n}=0}
\]

\end{tcolorbox}

\begin{tcolorbox}[title={Tips},colback=red!5!white
,colframe=red!75!black
,colbacktitle=yellow!50!red
,coltitle=red!25!black, fonttitle=\bfseries,subtitle style={boxrule=0.4pt, colback=yellow!50!red!25!white}]
判别级数的敛散性时，可先考察\(\lim\limits_{n \to \infty} \alpha_n \not= 0\)

\[
\text{如果}\begin{cases}
\lim\limits_{n \to \infty} \alpha_n \neq 0, & \text{级数发散} \\
\lim\limits_{n \to \infty} \alpha_n = 0, & \text{应进一步判断}
\end{cases}
\]

\end{tcolorbox}
\paragraph{绝对收敛和条件收敛}
\label{绝对收敛和条件收敛}
如果 \(\sum\limits_{n=1}^{\infty} \left| \alpha_n \right|\) 收敛，那么
\(\sum\limits_{n=1}^{\infty} \alpha_n\) 也收敛。

且不等式
\(\left| \sum\limits_{n=1}^{\infty} \alpha_n \right| \le \sum\limits_{n=1}^{\infty} \left| \alpha_n \right|\)
成立。

\begin{tcolorbox}[title={注意},colback=red!5!white
,colframe=red!75!black
,colbacktitle=yellow!50!red
,coltitle=red!25!black, fonttitle=\bfseries,subtitle style={boxrule=0.4pt, colback=yellow!50!red!25!white}]
\(\sum\limits_{n=1}^{\infty} \left| \alpha_n \right|\)
的各项都是非负的实数，应用正项级数的审敛法则判定。

\end{tcolorbox}

\textbf{定义}：如果 \(\sum\limits_{n=1}^{\infty}|\alpha_{n}|\) 收敛，那么称级数
\(\sum\limits_{n=1}^{\infty}\alpha_{n}\) 为\textbf{绝对收敛}。

非绝对收敛的收敛级数称为\textbf{条件收敛级数}。

\begin{tcolorbox}[title={说明},colback=red!5!white
,colframe=red!75!black
,colbacktitle=yellow!50!red
,coltitle=red!25!black, fonttitle=\bfseries,subtitle style={boxrule=0.4pt, colback=yellow!50!red!25!white}]
由 \(|\alpha_{n}|=\sqrt{a_{n}^{2}+b_{n}^{2}}\leq|a_{n}|+|b_{n}|\)，

所以 \(\sum\limits_{n=1}^{\infty}a_{n}\) 与
\(\sum\limits_{n=1}^{\infty}b_{n}\)
绝对收敛时，\(\sum\limits_{n=1}^{\infty}\alpha_{n}\) 也绝对收敛。

综上： \[
\boxed{\sum\limits_{n=1}^{\infty}\alpha_{n} \text{绝对收敛} \Leftrightarrow\sum\limits_{n=1}^{\infty}a_{n} \text{与} \sum\limits_{n=1}^{\infty}b_{n} \text{绝对收敛}}
\]

\end{tcolorbox}
\subsection{幂级数}
\label{幂级数}
\subsubsection{复变函数项级数、部分和}
\label{复变函数项级数部分和}
设 \(\{f_n(z)\}(n=1,2,\ldots)\) 为一复变函数序列，其中各项在区域 \(D\)
内有定义。表达式
\[
\sum_{n=1}^{\infty} f_n(z) = f_1(z) + f_2(z) + \cdots + f_n(z) + \cdots
\]
称为复变函数项级数，记作 \(\sum_{n=1}^{\infty}f_{n}(z)\)。
级数最前面 n 项的和
\[
s_{n}(z)=f_{1}(z)+f_{2}(z)+\cdots +f_{n}(z)
\]
称为这级数的\textbf{部分和}。
\subsubsection{和函数}
\label{和函数}
如果对于 \(D\) 内的某一点 \(z_0\)，若
\(\sum\limits_{n=1}^{\infty} f_n(z_0)\) 收敛，那末称级数
\(\sum\limits_{n=1}^{\infty} f_n(z)\) 在 \(z_0\) 收敛。其和记为
\(S(z_0)\)。

如果级数在 \(D\) 内处处收敛，那么它的和一定是 \(z\) 的一个函数
\(s(z)\)：

\[
s(z) = f_1(z) + f_2(z) + \cdots + f_n(z) + \cdots
\]

称为该级数在区域 \(D\) 上的\textbf{和函数}。
\subsubsection{幂级数概念}
\label{幂级数概念}
当 \(f_n(z) = c_{n-1}(z-a)^{n-1}\) 或 \(f_n(z) = c_{n-1}z^{n-1}\)
时，得到函数项级数的特殊情形：
\[
\sum_{n=0}^{\infty} c_n (z-a)^n = c_0 + c_1(z-a) + c_2(z-a)^2 + \cdots + c_n(z-a)^n + \cdots
\]
或
\[
\sum_{n=0}^{\infty} c_n z^n = c_0 + c_1 z + c_2 z^2 + \cdots + c_n z^n + \cdots
\]
这种级数称为幂级数。
\subsubsection{敛散性}
\label{敛散性}
\paragraph{收敛定理（Abel 定理）}
\label{收敛定理 abel 定理}
如果级数 \(\sum_{n=0}^{\infty} c_{n} z^{n}\) 在 \(z = z_{0} (\neq 0)\)
处收敛，那么对满足 \(|z| < |z_{0}|\) 的 \(z\)，级数必\textbf{绝对收敛}。

如果在 \(z = z_{0}\) 处级数发散，那么对满足 \(|z| > |z_{0}|\) 的
\(z\)，级数必\textbf{发散}。
\paragraph{收敛圆与收敛半径}
\label{收敛圆与收敛半径}
对于一个幕级数，其收敛半径的情况有三种：

\begin{enumerate}
\item 对所有的正实数都收敛。由阿贝尔定理知：\textbf{级数在复平面内处处绝对收敛}。
\item 对所有的正实数，除了 \(z=0\) 之外都发散。此时，\textbf{级数在复平面内除了原点之外处处
发散}。
\item 同时存在使得级数发散和收敛的正实数。这时，幂级数
\(\sum\limits_{n=0}^{\infty}c_nz^n\) 的收敛范围是一个以原点为中心的圆域。
\end{enumerate}

\begin{tcolorbox}[title={Example},colback=red!5!white
,colframe=red!75!black
,colbacktitle=yellow!50!red
,coltitle=red!25!black, fonttitle=\bfseries,subtitle style={boxrule=0.4pt, colback=yellow!50!red!25!white}]
对于 \(1+z+2^2z^2+\cdots+n^nz^n\) 这个级数，当
\(z \ne 0\) 时，通项是\textbf{不趋于零}的，级数发散。

\end{tcolorbox}
\paragraph{收敛半径的求法}
\label{收敛半径的求法}
\subparagraph{比值法}
\label{比值法}
\[
\boxed{
    \text{如果} \lim_{n \to \infty} \left|\frac{c_{n+1}}{c_n}\right| = \lambda \ne 0, \quad \text{则收敛半径} R=\frac{1}{\lambda}.
}
\]
\subparagraph{根植法}
\label{根植法}
\[
\boxed{
    \text{如果} \lim_{n \to \infty} \sqrt[n]{|c_n|} = \lambda \ne 0, \quad \text{则收敛半径} R = \frac{1}{\lambda}.
}
\]

\begin{tcolorbox}[title={注意},colback=red!5!white
,colframe=red!75!black
,colbacktitle=yellow!50!red
,coltitle=red!25!black, fonttitle=\bfseries,subtitle style={boxrule=0.4pt, colback=yellow!50!red!25!white}]
如果：
\begin{enumerate}
\item \(\lambda = 0\)，则级数 \(\sum\limits _{n=0}^{\infty}c_n z^n\) 在复平面内处处
收敛，即 \(\boxed{R=\infty}\)。
\item \(\lambda = \infty\)（极限不存在），则级数 \(\sum\limits _{n=0}^{\infty}c_n
   z^n\) 对于复平面内除 \(n=0\) 外的一切 \(z\) 均发散，也就是 \(\boxed{R=0}\)。
\end{enumerate}

\end{tcolorbox}
\subsubsection{幂级数的运算和性质}
\label{幂级数的运算和性质}
\paragraph{有理运算}
\label{有理运算}
设
\[
f(z)=\sum_{n=0}^{\infty} a_{n}z^{n},\quad R=r_{1},\qquad g(z)=\sum_{n=0}^{\infty} b_{n}z^{n},\quad R=r_{2}.
\]

则有：

\begin{align*}
f(z)\pm g(z) &=\sum_{n=0}^{\infty} a_{n}z^{n}\pm\sum_{n=0}^{\infty} b_{n}z^{n}\\
&= \sum_{n=0}^{\infty}(a_{n}\pm b_{n})z^{n}.\\
& (R=\min(r_{1},r_{2}), \quad |z| < R)\\
\\
f(z)\cdot g(z) &=\left(\sum_{n=0}^{\infty} a_{n}z^{n}\right)\cdot\left(\sum_{n=0}^{\infty} b_{n}z^{n}\right)\\
& =\sum_{n=0}^{\infty}\left(a_{n}b_{0}+a_{n-1}b_{1}+\cdots+a_{0} b_{n}\right)z^{n}.\\
& (R=\min(r_{1},r_{2}), \quad |z| < R)
\end{align*}
\paragraph{代换（复合）运算}
\label{代换复合运算}
\textbf{定理}： 如果当 \(|z|<r\)
时，\(f(z)=\sum\limits _{n=0}^{\infty} a_{n}z^{n}\)，又设在 \(|z|<R\) 内
\(g(z)\) 解析且满足 \(|g(z)|<r\)，那末当 \(|z|<R\)
时，\(f[g(z)]=\sum\limits_{n=0}^{\infty} a_{n}[g(z)]^{n}\)。

\begin{tcolorbox}[title={说明},colback=red!5!white
,colframe=red!75!black
,colbacktitle=yellow!50!red
,coltitle=red!25!black, fonttitle=\bfseries,subtitle style={boxrule=0.4pt, colback=yellow!50!red!25!white}]
此代换运算常应用于将函数展开成幂级数。

\end{tcolorbox}
\paragraph{复变幂函数在收敛圆内的性质}
\label{复变幂函数在收敛圆内的性质}
\textbf{定理} 设幂级数 \(\sum\limits_{n=0}^{\infty} c_{n}(z-a)^{n}\)
的收敛半径为 \(R\)，则

\begin{enumerate}
\item 它的和函数 \(f(z) = \sum\limits_{n=0}^{\infty} c_{n}(z-a)^{n}\)
是收敛圆 \(|z-a| < R\) 内的解析函数。
\item \(f(z)\) 在收敛圆 \(|z-a| < R\) 内的导数可将其幂级数逐项求导得到，即
\[
   f'(z) = \sum_{n=1}^{\infty} n c_{n}(z-a)^{n-1}.
   \]
\item 在收敛圆内可以逐项积分，也就是
\[
   \int\limits _c f(z)\mathrm{d}z = \sum _{n=0}^{\infty} c_n\int\limits _c (z-a)^n \mathrm{d}z, \quad c \subset \{ z | \left| z-a \right| < R \}.
   \]
或
\[
   \int _a^z f(\zeta)\mathrm{d}z = \sum _{n=0}^{\infty}\frac{c_n}{n+1}(z-a)^{n+1}.
   \]
\end{enumerate}

\begin{tcolorbox}[title={Note},colback=red!5!white
,colframe=red!75!black
,colbacktitle=yellow!50!red
,coltitle=red!25!black, fonttitle=\bfseries,subtitle style={boxrule=0.4pt, colback=yellow!50!red!25!white}]
简而言之，在收敛圆内，幂函数的\textbf{和函数解析}；幂函数可以\textbf{逐项求导}或者\textbf{逐项积分}。

\end{tcolorbox}
\subsection{泰勒级数}
\label{泰勒级数}
\subsubsection{泰勒定理}
\label{泰勒定理}
设 \(f(z)\) 在区域 \(D\) 内解析，\(z_0\) 为 \(D\)内的一点，\(d\) 为
\(z_0\) 到 \(D\) 的边界上各点的最短距离，则当 \(|z-z_0| < d\) 时，有
\(f(z)=\sum\limits_{n=0}^{\infty} c_n (z-z_0)^n\)成立。其中，
\(c_n = \dfrac{1}{n!}f^{(n)}(z_0),\quad n = 0,1,2,\cdots\)

\begin{tcolorbox}[title={说明},colback=red!5!white
,colframe=red!75!black
,colbacktitle=yellow!50!red
,coltitle=red!25!black, fonttitle=\bfseries,subtitle style={boxrule=0.4pt, colback=yellow!50!red!25!white}]
\begin{enumerate}
\item 复变函数展开为泰勒级数的条件与实函数展开为泰勒级数的条件有何不同：
\begin{quote}
复解析是一个很强的条件，它的效果要比实函数的无限阶可导还要强。因此对于实函数我们要考虑余项的极限是否为零，但是对于复变函数，只要复解析，就可以展开。
\end{quote}
\item 如果 \(f(z)\) 在 \(D\) 内有奇点，则 \(d\) 等于 \(z_{0}\) 到最近一个奇点 \(\alpha\) 之间的距离，即 \(d = |\alpha - z_0|\)。
\item 当 \(z_0 = 0\) 时，级数称为麦克劳林级数。
\item 任何解析函数在一点的泰勒级数是唯一的。
\end{enumerate}

\end{tcolorbox}
\subsubsection{函数展开为泰勒级数}
\label{函数展开为泰勒级数}
\paragraph{直接法}
\label{直接法}
用泰勒展开定理计算系数：
\[
c_n = \frac{1}{n!} f^{(n)}z_0,\quad n=0,1,2,\cdots
\]
例如：
\begin{align*}
    e^z &= 1 + z + \frac{z^2}{2!} + \cdots = \sum _{n=0}^{\infty} \frac{z^n}{n!}, \quad (R=\infty)
    \\
    \sin z &= z - \frac{z^3}{3!} + \frac{z^5}{5!} - \cdots + (-1)^n\frac{z^{2n+1}}{(2n+1)!} + \cdots,\quad (R=\infty)
    \\
    \cos z &= 1 - \frac{z^2}{2!} + \frac{z^4}{4!} - \cdots + (-1)^n\frac{z^{2n}}{(2n)!} + \cdots, \quad (R=\infty)
\end{align*}
\paragraph{间接法}
\label{间接法}
借助于一些已知函数的展开式，结合解析函数的性质，幂级数运算性质（逐项求导，积分等）和其它数学技巧（代换等），求函数的泰勒展开式。

以用间接法在 \(z=0\) 附近展开 \(\sin z\) 为例。

\begin{align*}
    \sin z &= \frac{1}{2i}(e^{iz}-e^{-iz})
    \\
    &= \frac{1}{2i}[\sum _{n=0}^{\infty}\frac{(iz)^n}{n!} - \sum _{n=0}^{\infty}\frac{(-iz)^n}{n!}]
    \\
    &= \sum _{n=0}^{\infty}(-1)^n\frac{z^{2n+1}}{(2n+1)!}
\end{align*}
\paragraph{一些常见的泰勒展开式}
\label{一些常见的泰勒展开式}
\begin{align*}
    e^{z} &= 1 + z + \frac{z^{2}}{2!} + \cdots + \frac{z^{n}}{n!} + \cdots
    \\
    &= \sum_{n=0}^{\infty} \frac{z^{n}}{n!}, \quad (|z| < \infty)
\\
    \frac{1}{1-z} &= 1 + z + z^{2} + \cdots + z^{n} + \cdots
    \\
    &= \sum_{n=0}^{\infty} z^{n}, \quad (|z| < 1)
\\
    \frac{1}{1+z} &= 1 - z + z^{2} - \cdots + (-1)^{n}z^{n} + \cdots
    \\
    &= \sum_{n=0}^{\infty} (-1)^{n}z^{n}, \quad (|z| < 1)
\\
    \sin z &= z - \frac{z^{3}}{3!} + \frac{z^{5}}{5!} - \cdots + (-1)^{n}\frac{z^{2n+1}}{(2n+1)!} + \cdots, \quad (|z| < \infty)
\\
    \cos z &= 1 - \frac{z^{2}}{2!} + \frac{z^{4}}{4!} - \cdots + (-1)^{n}\frac{z^{2n}}{(2n)!} + \cdots, \quad (|z| < \infty)
\\
    \ln(1+z) &= z - \frac{z^{2}}{2} + \frac{z^{3}}{3} - \cdots + (-1)^{n}\frac{z^{n+1}}{n+1} + \cdots \\
            &= \sum_{n=0}^{\infty} (-1)^{n} \frac{z^{n+1}}{n+1}, \quad (|z| < 1)
\\
    (1+z)^{\alpha} &= 1 + \alpha z + \frac{\alpha(\alpha-1)}{2!}z^{2} + \frac{\alpha(\alpha-1)(\alpha-2)}{3!}z^{3}
    \\
    &\quad + \cdots + \frac{\alpha(\alpha-1)\cdots(\alpha-n+1)}{n!}z^{n} + \cdots, \quad (|z| < 1)
\end{align*}
\subsection{洛朗级数}
\label{洛朗级数}
\subsubsection{双边幂级数}
\label{双边幂级数}
这种幂级数不仅包含正数次数的项，还包含负数次数的项。

\[
\sum_{n=-\infty}^{+\infty} c_n(z-z_0)^n = \sum_{n=1}^{\infty}
c_{-n}(z-z_0)^{-n} + \sum_{n=0}^{\infty} c_{n}(z-z_0)^{n}
\]

等号右边第一项是负幂项部分，是研究函数性质的主要部分，而第二项的正幂项部分是解析部分。

对于 \(\sum\limits_{n=1}^{\infty} c_{-n}(z-z_0)^{-n}\)，令
\(\zeta = (z-z_0)^{-1}\) 使其变为 \(\sum\limits_{n=1}^{\infty} c_{-n} \zeta^{n}\)
它的收敛半径记为 \(R\)，令 \(|\zeta| < R\) ，这时候收敛，记收敛域为
\(|z-z_0| > \dfrac{1}{R} = R_1\)。

对于 \(\sum\limits _{n=0}^{\infty} c_{n}(z-z_0)^{n}\)，它的收敛半径记为
\(R_2\)，收敛域 \(|z-z_0| < R_2\)。

接下来，若：

\begin{enumerate}
\item \(R_1 > R_2\)：两个收敛域无公共部分。
\item \(R_1 < R_2\)：两个收敛域有公共部分 \(R_1 < |z-z_0| < R_2\)。
\end{enumerate}

\begin{tcolorbox}[title={结论},colback=red!5!white
,colframe=red!75!black
,colbacktitle=yellow!50!red
,coltitle=red!25!black, fonttitle=\bfseries,subtitle style={boxrule=0.4pt, colback=yellow!50!red!25!white}]
双边幂级数的收敛域为圆环域
\(R_1 < |z-z_0| < R_2\)，在此圆环内和函数为解析函数。

常见的特殊圆环域：
\begin{enumerate}
\item \(0 < |z-z_0| < R_2\)，一个抠掉圆心的实心圆。
\item \(R_1 < |z-z_0| < \infty\)，从整个平面抠掉一个实心圆。
\item \(0 < |z-z_0| < \infty\)，从整个平面抠掉一个点。
\end{enumerate}

\end{tcolorbox}
\subsubsection{洛朗级数的概念}
\label{洛朗级数的概念}
设 \(f(z)\) 在圆环域 \(R_1 < |z - z_0| < R_2\) 内处处解析，那么 \(f(z)\)
在该圆环域 \(D\) 内可展开成\textbf{洛朗级数}：
\[
f(z) = \sum_{n=-\infty}^{\infty} c_n (z - z_0)^n
\]

其中，\textbf{洛朗系数} \(c_n\) 由积分定义：
\[
c_n = \frac{1}{2\pi i} \oint_C \frac{f(\zeta)}{(\zeta - z_0)^{n+1}} d\zeta \quad (n = 0, \pm 1, \cdots)
\]

这里，\(C\) 是圆环域内\textbf{绕 \(z_0\) 的任一正向简单闭曲线}。
\subsubsection{函数的洛朗展开}
\label{函数的洛朗展开}
\paragraph{直接法}
\label{直接法-1}
直接利用上面的定理公式计算洛朗系数，然后写出洛朗展开式。但是，计算往往很麻烦。
\paragraph{间接法}
\label{间接法-1}
根据正负幂项组成的级数的唯一性，可以用代数运算，代换，求导，积分等方法去展开。
